\documentclass[11pt,a4paper]{article}

\usepackage[margin=2.5cm]{geometry}
\usepackage{amsmath, amssymb}
\usepackage{hyperref}
\usepackage{listings}
\usepackage{xcolor}
\usepackage{booktabs}
\usepackage{parskip}

\lstset{
  basicstyle=\ttfamily\small,
  keywordstyle=\color{blue},
  commentstyle=\color{gray},
  stringstyle=\color{teal},
  breaklines=true,
  frame=single,
  backgroundcolor=\color{gray!8},
  columns=fullflexible,
}

\title{Merge-and-Shrink Algorithm:\\Implementation Overview}
\author{ms-axioms / Fast Downward}
\date{}

\begin{document}

\maketitle
\tableofcontents
\newpage

% ─────────────────────────────────────────────────────────────────────────────
\section{Overview}

The merge-and-shrink (M\&S) framework builds an admissible heuristic for
cost-optimal classical planning by constructing a compressed abstract
representation of the planning task.  Starting from one atomic transition
system per planning variable, the algorithm iteratively:

\begin{enumerate}
  \item \textbf{Shrinks} two chosen systems to keep their sizes manageable, then
  \item \textbf{Merges} them into a single system via synchronized product.
\end{enumerate}

The resulting abstract system encodes goal distances that lower-bound the true
cost to the goal, yielding an admissible (and, under standard conditions,
consistent) heuristic.

% ─────────────────────────────────────────────────────────────────────────────
\section{Key Source Files}

\begin{center}
\begin{tabular}{ll}
\toprule
\textbf{File} & \textbf{Role} \\
\midrule
\texttt{merge\_and\_shrink\_algorithm.cc} & Algorithm orchestration \\
\texttt{factored\_transition\_system.cc} & Manages the collection of factors \\
\texttt{transition\_system.cc} & Individual abstract state space \\
\texttt{shrink\_bisimulation.cc} & Primary shrink strategy \\
\texttt{label\_reduction.cc} & Merges redundant action labels \\
\texttt{merge\_and\_shrink\_heuristic.cc} & Planner-facing heuristic interface \\
\texttt{merge\_and\_shrink\_representation.cc} & Abstract state lookup tables \\
\texttt{distances.cc} & BFS for init/goal distances \\
\bottomrule
\end{tabular}
\end{center}

% ─────────────────────────────────────────────────────────────────────────────
\section{Algorithm Flow}

The entry point is \texttt{MergeAndShrinkAlgorithm::build\_factored\_transition\_system()}.
After creating one atomic transition system per variable and optionally pruning
unreachable or dead-end states, it runs the main loop:

\begin{lstlisting}[caption={Main loop pseudocode}]
while |active factors| > 1:
    (i1, i2)  <- merge_strategy.get_next()          # 1. select pair
    label_reduction.reduce(i1, i2)                   # 2. optional, before shrink
    shrink_before_merge(TS[i1], TS[i2])              # 3. shrink both
    label_reduction.reduce(i1, i2)                   # 4. optional, before merge
    i  <- fts.merge(i1, i2)                          # 5. synchronized product
    prune(TS[i])                                     # 6. remove dead states
    if not fts.is_factor_solvable(i): break          # 7. early termination
\end{lstlisting}

% ─────────────────────────────────────────────────────────────────────────────
\section{Merge Strategies}

The merge strategy decides \emph{which} pair of factors to combine next.
Three families are implemented.

\subsection{Precomputed (\texttt{merge\_strategy\_precomputed.cc})}

A \texttt{MergeTree} --- a full binary tree over variables --- is constructed
before the loop and consulted at each step.  The next merge is always the
leftmost sibling-leaf pair in the remaining tree.  The canonical choice is a
\emph{linear} tree ordered by reverse level in the causal graph.

\subsection{Stateless (\texttt{merge\_strategy\_stateless.cc})}

Pairs are chosen dynamically using a \texttt{MergeSelector} that scores every
candidate pair and keeps only those that maximise the score.  Available scoring
functions include:

\begin{itemize}
  \item \textbf{Goal relevance} --- prefer factors that share goal variables.
  \item \textbf{DFP} --- the classic Dräger--Finkbeiner--Podelski criterion.
  \item \textbf{MIASM} --- maximise abstract state space size reduction.
  \item \textbf{Total order} --- tie-breaking by a fixed variable ordering.
\end{itemize}

\subsection{SCC-based (\texttt{merge\_strategy\_sccs.cc})}

Merges are first restricted to strongly connected components of the causal
graph, which are processed in topological order.  A configurable fallback
strategy handles merges across SCCs.

% ─────────────────────────────────────────────────────────────────────────────
\section{Shrink Strategies}

Before merging, each factor is reduced so the product stays within the
\texttt{max\_states} budget.  All strategies compute a
\emph{state equivalence relation} --- a partition of states into equivalence
classes --- which is then applied to the transition system.

\subsection{Size Computation (\texttt{utils.cc})}

Let $s_1, s_2$ be the current sizes and $M$ the hard limit.  Target sizes
$(t_1, t_2)$ satisfy:

\begin{align}
  t_i &\leq \texttt{max\_states\_before\_merge} \\
  t_1 \cdot t_2 &\leq M
\end{align}

If both $s_i > \sqrt{M}$, both targets are set to $\lfloor\sqrt{M}\rfloor$.
Otherwise the larger one is reduced to $\lfloor M / s_{\min} \rfloor$.

\subsection{Bisimulation (\texttt{shrink\_bisimulation.cc})}

The strongest strategy.  States $s, s'$ are bisimilar if they have identical
goal-distance $h$ and, for every abstract transition $(s, \ell, t)$, there
exists a transition $(s', \ell, t')$ with $t \sim t'$, and vice versa.

\paragraph{Algorithm.}
\begin{enumerate}
  \item \textbf{Initialise groups.}  States are partitioned by $h$-value:
        goal states form group~0; states with $h = i > 0$ form group~$i$;
        dead-end states ($h = \infty$) get a dedicated group.
  \item \textbf{Compute signatures.}  Each state $s$ gets a signature
        \[
          \sigma(s) = \bigl(h(s),\; \text{group}(s),\;
            \{(\ell, \text{group}(t)) : (s,\ell,t) \in T\}\bigr).
        \]
  \item \textbf{Refine.}  States with identical signatures are merged into one
        group.  Repeat until a fixpoint is reached.
\end{enumerate}

With \texttt{greedy=true}, only $h$-optimal transitions are included in the
signatures, preserving at least one optimal plan while allowing more aggressive
merging.

\subsection{Bucket-based (\texttt{shrink\_fh.cc}, \texttt{shrink\_random.cc})}

Simpler strategies that group states into buckets by $(f, h)$ values or
randomly, then merge states within each bucket to meet the target size.

% ─────────────────────────────────────────────────────────────────────────────
\section{Label Reduction (\texttt{label\_reduction.cc})}

Action labels that induce \emph{identical transitions} across all relevant
factors can be merged without changing the abstract distances.  This reduces
the branching factor of every transition system at no precision cost.

Three application timings are supported:

\begin{itemize}
  \item \texttt{TWO\_TRANSITION\_SYSTEMS} --- only for the next merge pair.
  \item \texttt{ALL\_TRANSITION\_SYSTEMS} --- across all active factors.
  \item \texttt{ALL\_TRANSITION\_SYSTEMS\_WITH\_FIXPOINT} --- iteratively until
        no further reduction is possible.
\end{itemize}

% ─────────────────────────────────────────────────────────────────────────────
\section{Distance Computation (\texttt{distances.cc})}

After each merge, two breadth-first searches are run on the new factor:

\begin{itemize}
  \item \textbf{Forward BFS} from the initial abstract state computes
        $\mathit{init\_distance}(s)$ for every state $s$.
  \item \textbf{Backward BFS} from all goal states computes
        $\mathit{goal\_distance}(s)$ (i.e.\ $h(s)$).
\end{itemize}

States with $\mathit{init\_distance} = \infty$ are unreachable;
states with $\mathit{goal\_distance} = \infty$ are dead ends.
Both can be pruned to keep the factor compact.

% ─────────────────────────────────────────────────────────────────────────────
\section{Abstract State Representation (\texttt{merge\_and\_shrink\_representation.cc})}

After the algorithm terminates, the heuristic is stored as a tree of lookup
tables mirroring the merge tree:

\begin{itemize}
  \item \textbf{Leaf nodes} (\texttt{MergeAndShrinkRepresentationLeaf}) map
        concrete variable values to abstract states via a flat array.
  \item \textbf{Internal nodes} (\texttt{MergeAndShrinkRepresentationMerge})
        combine two children through a 2-D table indexed by the children's
        abstract state indices.
\end{itemize}

Once the algorithm finishes, \texttt{set\_distances()} replaces every abstract
state index with its goal distance, so that evaluating the heuristic for a
concrete state is a pure table lookup with no search.

% ─────────────────────────────────────────────────────────────────────────────
\section{Configuration Example}

The recommended configuration uses SCC-based merging, non-greedy bisimulation
shrinking, and label reduction before shrinking:

\begin{lstlisting}[caption={Recommended M\&S configuration}]
astar(merge_and_shrink(
    merge_strategy=merge_sccs(
        order_of_sccs=topological,
        merge_selector=score_based_filtering(
            scoring_functions=[goal_relevance(), dfp(), total_order()])),
    shrink_strategy=shrink_bisimulation(greedy=false),
    label_reduction=exact(before_shrinking=true, before_merging=false),
    max_states=50000,
    threshold_before_merge=1))
\end{lstlisting}

To run this on the Miconic benchmark included in the repository:

\begin{lstlisting}[language=bash, caption={Running on Miconic s1-0}]
./fast-downward.py \
    misc/tests/benchmarks/miconic/domain.pddl \
    misc/tests/benchmarks/miconic/s1-0.pddl \
    --search "astar(merge_and_shrink(
        merge_strategy=merge_sccs(
            order_of_sccs=topological,
            merge_selector=score_based_filtering(
                scoring_functions=[goal_relevance(),dfp(),total_order()])),
        shrink_strategy=shrink_bisimulation(greedy=false),
        label_reduction=exact(before_shrinking=true,before_merging=false),
        max_states=50000,
        threshold_before_merge=1))"
\end{lstlisting}

% ─────────────────────────────────────────────────────────────────────────────
\section{Algorithm Properties}

\begin{center}
\begin{tabular}{ll}
\toprule
\textbf{Property} & \textbf{Value} \\
\midrule
Admissibility & Guaranteed (no pruning of reachable optimal states) \\
Consistency & Guaranteed under standard conditions \\
Completeness & May terminate early; uses partial abstractions as heuristics \\
Input & STRIPS/ADL tasks without axioms \\
Size control & \texttt{max\_states}, \texttt{max\_states\_before\_merge} \\
Time control & \texttt{main\_loop\_max\_time} \\
\bottomrule
\end{tabular}
\end{center}

\end{document}
